\documentclass[11pt]{article}

% ---------------------------------------------------------------------------
% shifty: the theory and engineering behind a formalism-first SHACL engine.
%
% WORKING DRAFT. Bullet-point + math density on purpose; prose is intentionally
% thin so the technical content (algebra, transformations, SPARQL lowering,
% repair) is captured before it is polished into a publication.
% ---------------------------------------------------------------------------

\usepackage[margin=1in]{geometry}
\usepackage{amsmath,amssymb,amsthm}
\usepackage{mathtools}
\usepackage{booktabs}
\usepackage{array}
\usepackage{enumitem}
\usepackage{listings}
\usepackage{xcolor}
\usepackage{microtype}
\usepackage[hidelinks]{hyperref}
\usepackage{tikz}
\usetikzlibrary{arrows.meta,positioning}

% --- notation -------------------------------------------------------------
\newcommand{\sem}[1]{[\![#1]\!]}                 % denotation brackets
\newcommand{\semG}[1]{\sem{#1}^{G}}              % denotation over graph G
\newcommand{\Term}{\mathit{Term}}
\newcommand{\Node}{N}
\newcommand{\Val}{V}
\newcommand{\Pred}{P}
\newcommand{\Key}{K}
\newcommand{\conf}{\models}                       % G,v |= phi
\newcommand{\id}{\mathsf{id}}
\newcommand{\test}{\mathsf{test}}
\newcommand{\closedc}{\mathsf{closed}}
\newcommand{\eqc}{\mathsf{eq}}
\newcommand{\disjc}{\mathsf{disj}}
\newcommand{\kind}{\mathsf{kind}}
\newcommand{\sel}{\mathit{sel}}
\newcommand{\lfp}{\mathrm{lfp}}
\newcommand{\gfp}{\mathrm{gfp}}
\newcommand{\Star}{{*}}
\newcommand{\Repair}{\mathsf{repair}}
\newcommand{\Break}{\mathsf{break}}
\newcommand{\Build}{\mathsf{build}}
\newcommand{\code}[1]{\texttt{#1}}

\theoremstyle{definition}
\newtheorem{definition}{Definition}[section]
\newtheorem{law}{Law}[section]
\newtheorem{principle}{Principle}[section]

% --- listings -------------------------------------------------------------
\definecolor{kw}{rgb}{0.30,0.20,0.65}
\definecolor{cm}{rgb}{0.35,0.45,0.35}
\definecolor{st}{rgb}{0.60,0.25,0.20}
\lstdefinelanguage{rustlite}{
  morekeywords={enum,struct,fn,impl,pub,let,mut,match,Box,Vec,Option,Some,None,
    Term,Path,Shape,self,for,in,if,else,return,u32,u64,bool,HashMap,HashSet,Rc,
    RefCell,BTreeMap,BTreeSet,where,trait,type,as,move,crate},
  sensitive=true,
  morecomment=[l]{//},
  morecomment=[s]{/*}{*/},
  morestring=[b]",
}
\lstset{
  basicstyle=\ttfamily\small,
  keywordstyle=\color{kw}\bfseries,
  commentstyle=\color{cm}\itshape,
  stringstyle=\color{st},
  columns=fullflexible,
  breaklines=true,
  showstringspaces=false,
  frame=single,
  framesep=4pt,
  rulecolor=\color{black!25},
  xleftmargin=4pt,
}
\lstset{language=rustlite}

\newcommand{\tagdone}{{\small\textsc{[impl]}}}
\newcommand{\tagdesign}{{\small\textsc{[design]}}}

% --- implementation-status callouts ---------------------------------------
% Inline marker and a framed block for material that is DESIGNED but NOT YET
% in the current crates. Kept (not deleted) so we can decide later what to cut.
\definecolor{nired}{rgb}{0.60,0.05,0.05}
\newcommand{\NI}[1]{{\color{nired}\textbf{[NOT IMPLEMENTED --- design only:} #1\textbf{]}}}
\newcommand{\PART}[1]{{\color{nired}\textbf{[PARTIAL:} #1\textbf{]}}}
\newcommand{\nibox}[1]{%
  \par\smallskip\noindent
  \fcolorbox{nired}{nired!6}{%
    \parbox{\dimexpr\linewidth-2\fboxsep-2\fboxrule}{%
      \small\textbf{\color{nired}NOT IMPLEMENTED (design only --- candidate for removal later).}\ #1}}%
  \par\smallskip}

\title{\textbf{shifty}: A Formalism-First SHACL / SHACL-AF Engine\\
\large Algebraic IR, Semantics-Preserving Optimization, Native SPARQL Lowering,\\ and Symbolic Repair}
\author{Gabe Fierro \and (working notes)}
\date{Draft --- \today}

\begin{document}
\maketitle

\begin{abstract}
\noindent
\textbf{shifty} is a SHACL Core + SHACL-AF (Advanced Features) validation and
inference engine that is \emph{formalism-first}: every shapes graph is lowered to
a small algebraic intermediate representation (IR) --- a path algebra $\pi$ and a
shape grammar $\varphi$ adapted from the \emph{Common Foundations for SHACL,
ShEx, and PG-Schema} treatment (Ahmetaj et al., arXiv:2502.01295), specialized to
RDF --- before anything is evaluated. The same IR drives validation, inference,
static analysis, physical planning, native SPARQL execution, and a symbolic
\emph{repair} subsystem. This document records, in bullet-and-math form, (i) the
core formalism and the deliberate deviations needed to reach W3C coverage; (ii)
the pinned recursion semantics (stratified, greatest-fixpoint validation /
least-fixpoint inference) and the \emph{semantic} polarity analysis it rests on;
(iii) the catalog of semantics-preserving algebraic transformations; (iv) the
logical-to-physical planner and its cost model; (v) the native SPARQL lowering
over an immutable dictionary-encoded dataset, which is the principal performance
mechanism; and (vi) the repair calculus, which expresses graph repair as
\emph{abduction over $\varphi$} with a witness/satisfaction-trace duality and
multiple solver back-ends (monomorphism, enumeration, ASP, LLM). The emphasis is
on the engineering and the inventions that make the engine fast.
\end{abstract}

\tableofcontents

% ===========================================================================
\section{Overview and design thesis}
% ===========================================================================

\paragraph{The thesis.} W3C SHACL is specified as a large vocabulary of
constraint components with overlapping semantics. shifty's bet is that almost all
of that vocabulary collapses into a tiny algebra, and that \emph{optimizing one
small algebra} beats optimizing dozens of vocabulary terms.

\begin{itemize}[leftmargin=1.4em]
  \item \textbf{One reference semantics, many engines.} A naive denotational
  evaluator (the \emph{oracle}) defines truth. Every optimized engine must
  produce identical conformance results; this is enforced continuously against
  the W3C test suite and as property tests of the form
  $\mathsf{validate}(\mathsf{normalize}(S),G) \equiv \mathsf{validate}(S,G)$.
  \item \textbf{Lower early, optimize late.} All W3C vocabulary sugar is
  destructured into the core IR \emph{at parse time}. Optimizers and planners
  never see SHACL vocabulary --- only $\pi$/$\varphi$.
  \item \textbf{Fix semantics before optimizing.} The recursion / fixpoint
  semantics (\S\ref{sec:recursion}) is pinned before any rewrite that depends on
  it, so soundness of every transformation has a single referent.
\end{itemize}

\paragraph{Pipeline.} The stages are independently inspectable
(\code{shifty inspect --stage rdf|algebra|normalized|strata|plan|capability}):
\[
\underbrace{\text{RDF}}_{\text{Turtle}}
\;\xrightarrow{\text{parse + lower}}\;
\underbrace{(\pi,\varphi)\text{ IR}}_{\text{Schema arena}}
\;\xrightarrow{\text{normalize + CSE}}\;
\underbrace{\text{normal form}}_{\text{stratified}}
\;\xrightarrow{\text{plan}}\;
\underbrace{\text{physical plan}}_{\text{cost-ordered}}
\;\xrightarrow{\text{execute}}\;
\underbrace{\text{report / } \Delta G}_{\text{validate / infer / repair}}
\]

\paragraph{Implementation-status convention.} This is a working draft that
captures both what is \emph{built} and what is \emph{designed}. Material that is
specified but \textbf{not yet implemented in the current crates} is flagged
inline with \NI{\dots} or in a framed \emph{NOT IMPLEMENTED} box, and partial
implementations with \PART{\dots}. These markers are deliberate: they let us
decide later whether to keep the design material in a publication or cut it. Code
sketches without a marker reflect types/functions that exist in the tree.

\paragraph{Crate decomposition.}
\begin{center}
\begin{tabular}{@{}ll@{}}
\toprule
crate & role \\
\midrule
\code{shifty-algebra} & core IR: path algebra $\pi$, shape grammar $\varphi$, value types, schema arena, rendering \\
\code{shifty-parse}   & Turtle/RDF $\rightarrow$ IR lowering (all Core + AF sugar) \\
\code{shifty-opt}     & normalization/CSE, stratification, planning, native SPARQL lowering \\
\code{shifty-engine}  & validation, AF inference, SPARQL executor, repair (witness/synthesize/gate) \\
\code{shifty-cli}     & \code{shifty} binary; per-stage inspection \\
\code{pyshifty}       & PyO3 bindings (\code{validate}, \code{infer}, \code{RepairSession}) \\
\bottomrule
\end{tabular}
\end{center}

% ===========================================================================
\section{The algebraic core: $\pi$ and $\varphi$}\label{sec:core}
% ===========================================================================

The reference is the SHACL fragment of the \emph{Common Foundations} paper,
\textbf{specialized to RDF} and stripped of property-graph / ShEx / PG-Schema
material.

\subsection{Data model and the RDF re-unification (deviations D0, D1)}

\begin{itemize}[leftmargin=1.4em]
  \item The paper works over \emph{common graphs} $G=(E,\rho)$ with
  node-to-node edges $E\subseteq \Node\times \Pred\times \Node$ and a partial
  value map $\rho:\Node\times \Key \rightharpoonup \Val$. It deliberately keeps
  predicates $\Pred$ (edges) and keys $\Key$ (properties) \emph{distinct}, so
  paths navigate only node$\to$node and never ``through'' a value.
  \item \textbf{D0 --- re-unify $\Pred$ and $\Key$.} For RDF this wall is
  artificial; a common graph is just a set of triples. We collapse:
  \[
  \Term = \mathrm{IRI}\cup\mathrm{Blank}\cup\mathrm{Literal},\quad
  \Node=\mathrm{IRI}\cup\mathrm{Blank},\quad
  \Val=\mathrm{Literal},\quad
  \Pred=\Key=\mathrm{IRI},
  \]
  \[
  G\subseteq \Node\times\mathrm{IRI}\times\Term.
  \]
  Consequently a path may traverse \emph{any} triple regardless of object kind,
  so path denotation is a relation over $\Term\times\Term$, not $\Node\times\Node$.
  \item \textbf{D1 --- node identity is visible.} The paper hides node identity
  (which is \emph{why} it omits \code{sh:nodeKind}, class info, and node-constant
  comparison). Making identity visible re-enables those (extensions K1, C1, V1
  in \S\ref{sec:gap}).
\end{itemize}

\paragraph{Value types $T$.} A value type $\tau\in T$ denotes a decidable set
$\sem{\tau}\subseteq \Val$, with top type $\mathsf{any}$
($\sem{\mathsf{any}}=\Val$). This single abstraction absorbs datatypes, numeric
bounds, length bounds, regex patterns, and (as an extension) language membership:
\begin{lstlisting}
enum ValueType {
    Any,
    Datatype(IriBuf),                                   // sh:datatype
    NumericRange { lo: Option<Bound>, hi: Option<Bound> }, // sh:min/maxIn/Exclusive
    Length { min: Option<u64>, max: Option<u64> },      // sh:min/maxLength
    Pattern { regex: String, flags: String },           // sh:pattern (+ flags)
    LangIn(Vec<String>),                                 // sh:languageIn  (ext L1)
    And(Vec<ValueType>),                                 // intersection of facets
}
\end{lstlisting}

\subsection{Path expressions $\pi$ (Kleene algebra with converse)}

\begin{definition}[Path syntax]
$\pi ::= \id \mid q \mid \pi^{-} \mid \pi\cdot\pi' \mid \pi\cup\pi' \mid \pi^{\Star}$, for $q\in\Pred$.
\end{definition}

\noindent Denotation $\semG{\pi}\subseteq \Term\times\Term$:
\[
\begin{array}{ll}
\semG{\id} = \{(v,v)\} &
\semG{q} = \{(v,u)\mid (v,q,u)\in G\}\\[2pt]
\semG{\pi^{-}} = \{(v,u)\mid (u,v)\in\semG{\pi}\} &
\semG{\pi\cup\pi'} = \semG{\pi}\cup\semG{\pi'}\\[2pt]
\semG{\pi\cdot\pi'} = \{(v,w)\mid \exists u.\,(v,u)\in\semG{\pi}\wedge(u,w)\in\semG{\pi'}\} &
\semG{\pi^{\Star}} = \text{reflexive-transitive closure of } \semG{\pi}
\end{array}
\]
\begin{itemize}[leftmargin=1.4em]
  \item $\pi^{+}$ and $\pi^{?}$ are \emph{not} IR constructors: they are
  normalized away in the parser as $\pi^{+}=\pi\cdot\pi^{\Star}$ and
  $\pi^{?}=\pi\cup\id$ (extension P1).
  \item The carrier is a \emph{Kleene algebra with converse}; the laws this
  yields are exploited in normalization (\S\ref{sec:norm}).
\end{itemize}

\subsection{Shape grammar $\varphi$ (Boolean algebra over typed atoms)}

\begin{definition}[Shape syntax]
\[
\varphi ::= \top \mid \test(c) \mid \test(\tau) \mid \closedc(Q) \mid \eqc(\pi,p)
\mid \disjc(\pi,p) \mid \neg\varphi \mid \varphi\wedge\varphi' \mid
\varphi\vee\varphi' \mid \exists^{\geq n}\pi.\varphi \mid \exists^{\leq n}\pi.\varphi
\]
with $c\in\Val$, $\tau\in T$, $Q\subseteq_{\mathrm{fin}}\Pred\cup\Key$, $p\in\Pred$, $n\in\mathbb{N}$.
\end{definition}

\noindent Satisfaction $G,v\conf\varphi$ for a focus $v\in\Node\cup\Val$:
\[
\begin{array}{ll}
G,v\conf\top & \text{always}\\
G,v\conf\test(c) & v=c\\
G,v\conf\test(\tau) & v\in\sem{\tau}\\
G,v\conf\closedc(Q) & \forall p\in(\Pred\cup\Key)\setminus Q:\ \neg(G,v\conf\exists^{\geq1}p.\top)\\
G,v\conf\eqc(\pi,p) & \{u\mid(v,u)\in\semG{\pi}\} = \{u\mid(v,u)\in\semG{p}\}\\
G,v\conf\disjc(\pi,p) & \{u\mid(v,u)\in\semG{\pi}\} \cap \{u\mid(v,u)\in\semG{p}\} = \emptyset\\
G,v\conf\exists^{\geq n}\pi.\varphi & \#\{u\mid(v,u)\in\semG{\pi}\wedge G,u\conf\varphi\}\geq n\\
G,v\conf\exists^{\leq n}\pi.\varphi & \#\{u\mid(v,u)\in\semG{\pi}\wedge G,u\conf\varphi\}\leq n
\end{array}
\]
Derived sugar (kept in the parser, normalized away):
$\exists\pi.\varphi \triangleq \exists^{\geq1}\pi.\varphi$ and
$\forall\pi.\varphi \triangleq \exists^{\leq0}\pi.\neg\varphi$.

\paragraph{Fusion of the two count rows.} The IR fuses $\exists^{\geq n}$ and
$\exists^{\leq n}$ into a single \code{Count} constructor, because real SHACL
\emph{emits them as a pair} on a shared path and qualifier
(\code{sh:minCount}+\code{sh:maxCount}, or
\code{sh:qualifiedMin/MaxCount}); keeping them apart would force the optimizer to
re-pair them. The two-row semantics is recovered by treating $\min$ and $\max$
independently.
\begin{lstlisting}
enum Shape {
    Top,
    TestConst(Term),          // test(c), widened to Term (ext V1)
    TestType(ValueType),      // test(tau)
    Closed(BTreeSet<IriBuf>), // closed(Q): Q is the *allowed* predicate set
    Eq(Path, IriBuf), Disj(Path, IriBuf),
    Not(Box<Shape>), And(Vec<Shape>), Or(Vec<Shape>),
    Count { path: Path, min: Option<u64>, max: Option<u64>, qualifier: Box<Shape> },
    // --- extensions beyond the paper core (see gap analysis) ---
    TestKind(NodeKindSet),    // sh:nodeKind         (ext K1)
    Lt(Path, IriBuf), Le(Path, IriBuf),  // sh:lessThan(OrEquals) (ext O1)
    UniqueLang(Path),         // sh:uniqueLang       (ext L1)
}
\end{lstlisting}

\subsection{Selectors and schemas}

\begin{definition}[Selectors]
$\sel ::= \exists q.\top \mid \exists q^{-}.\top \mid \test(c)$. A \emph{schema}
$S$ is a finite set of pairs $(\sel,\varphi)$, and
\[
G\models S \iff \forall v\in\Node\cup\Val.\ \forall(\sel,\varphi)\in S.\ (G,v\conf\sel)\Rightarrow(G,v\conf\varphi).
\]
\end{definition}
\begin{itemize}[leftmargin=1.4em]
  \item The paper's selectors are unary and very restricted. shifty adds two:
  \code{HasPath(Path, Shape)} ($\exists^{\geq1}\pi.\varphi$, used for
  \code{sh:targetClass} / implicit class targets) and \code{Sparql(\dots)} (the
  AF SPARQL-target escape hatch).
  \item Shape references are resolved into a shared \emph{arena}; cyclic
  references are allowed (semantics in \S\ref{sec:recursion}).
\end{itemize}

\subsection{What the core buys}
\begin{itemize}[leftmargin=1.4em]
  \item A \textbf{closed denotational} definition of $G\models S$ --- decidable
  by structural recursion on $\varphi$ plus relational evaluation of $\pi$. This
  is the oracle.
  \item \textbf{Algebraic laws} to exploit: $\pi$ is a Kleene algebra with
  converse; $\varphi$ is a Boolean algebra over the
  $\test/\code{Count}/\eqc/\disjc/\closedc$ atoms.
  \item A \textbf{single counting primitive} \code{Count} subsuming cardinality,
  qualified cardinality, \code{sh:node}/\code{sh:property} nesting, and the
  $\forall/\exists$ sugar --- so the planner optimizes one construct, not a dozen.
\end{itemize}

% ===========================================================================
\section{Lowering W3C SHACL/SHACL-AF into the core}\label{sec:gap}
% ===========================================================================

The core is elegant but narrow (common graphs, validation only, non-recursive).
Reaching W3C \textbf{SHACL Core} + \textbf{SHACL-AF} requires a precise
``holes register.'' Each item is \emph{covered} ($\checkmark$, expressible
as-is or via documented sugar), an \emph{extension} ($\triangle$, a new IR
constructor that is still algebraic/optimizable), or an \emph{escape-hatch}
($\square$, opaque SPARQL/JS: compile a subset, fall back otherwise).

\subsection{Cleanly covered ($\checkmark$): sugar over the core}
\begin{center}
\small
\begin{tabular}{@{}ll@{}}
\toprule
W3C component & lowering \\
\midrule
\code{sh:datatype} / \code{sh:min/maxLength} / \code{sh:pattern} & $\test(\tau)$ facets \\
\code{sh:min/maxInclusive}, \code{sh:min/maxExclusive} & $\test(\tau)$ \code{NumericRange} \\
\code{sh:minCount} $n$ / \code{sh:maxCount} $n$ on $\pi$ & $\exists^{\geq n}\pi.\top$ / $\exists^{\leq n}\pi.\top$ \\
\code{sh:qualifiedValueShape}+\code{qualifiedMin/MaxCount} & $\exists^{\geq n}\pi.\varphi$ / $\exists^{\leq n}\pi.\varphi$ \\
\code{sh:node} $S$ (in a property shape) & $\forall\pi.\varphi_S = \exists^{\leq0}\pi.\neg\varphi_S$ \\
\code{sh:and}/\code{sh:or}/\code{sh:not} & $\wedge$/$\vee$/$\neg$ \\
\code{sh:xone} & $\bigvee_i(\varphi_i\wedge\bigwedge_{j\neq i}\neg\varphi_j)$ \\
\code{sh:closed} (+\code{sh:ignoredProperties}) & $\closedc(Q)$, $Q=\{\text{used paths}\}\cup\text{ignored}$ \\
\code{sh:equals} / \code{sh:disjoint} & $\eqc(\pi,p)$ / $\disjc(\pi,p)$ \\
\code{sh:hasValue} $c$ / \code{sh:in} $(c_1\dots c_k)$ & $\exists^{\geq1}\pi.\test(c)$ / $\exists^{\geq1}\pi.\bigvee_i\test(c_i)$ \\
\code{sh:targetSubjectsOf}/\code{ObjectsOf} $p$ & $\exists p.\top$ / $\exists p^{-}.\top$ \\
\bottomrule
\end{tabular}
\end{center}

\subsection{Extensions ($\triangle$): new but still algebraic}
\begin{itemize}[leftmargin=1.4em]
  \item \textbf{K1} \code{sh:nodeKind} $\Rightarrow$ \code{TestKind(NodeKindSet)},
  a primitive on $\kind(v)$ over $\{$IRI, Blank, Literal$\}$ and unions.
  \item \textbf{C1} \code{sh:class}/\code{sh:targetClass} as \emph{pure sugar
  over the path algebra} --- no class-specific machinery:
  \[
  \text{\code{sh:class} } C \ \equiv\ \exists^{\geq1}\,(\code{rdf:type}\cdot\code{rdfs:subClassOf}^{\Star}).\test_{\mathrm{node}}(C),
  \]
  and \code{sh:targetClass} is the same expression as a \code{HasPath} selector.
  The class-hierarchy path is configurable.
  \item \textbf{V1} widen $\test(c)$ from $c\in\Val$ to $c\in\Term$
  (\code{TestConst(Term)}) --- semantics unchanged, larger constant domain;
  enables IRI-valued \code{hasValue}/\code{in}/\code{targetNode}.
  \item \textbf{O1} \code{sh:lessThan(OrEquals)} $\Rightarrow$
  \code{Lt}/\code{Le} over the SPARQL term ordering.
  \item \textbf{L1} \code{sh:languageIn} $\Rightarrow$ $\test(\tau)$ \code{LangIn};
  \code{sh:uniqueLang} $\Rightarrow$ \code{UniqueLang}$(\pi)$ (counting over
  language tags).
  \item \textbf{P1} $\pi^{?}$/$\pi^{+}$ are parser-side sugar (no IR change).
\end{itemize}

\subsection{SHACL-AF: two whole subsystems}
SHACL-AF is absent from the paper and is the bulk of the new work.
\begin{itemize}[leftmargin=1.4em]
  \item \textbf{AF-R Rules.} A separate \emph{rule algebra}:
  \code{Rule \{ selector, conditions: Vec<Shape>, head: RuleHead, order \}}.
  \begin{itemize}
    \item \code{TripleRule} head $=(\text{subjectExpr},p,\text{objectExpr})$ built
    from node expressions;
    \item \code{SPARQLRule} head $=$ \code{CONSTRUCT} ($\square$).
  \end{itemize}
  Semantics: add inferred triples, iterate to a \emph{fixpoint} respecting
  \code{sh:order} and \code{sh:condition} (\S\ref{sec:infer}).
  \item \textbf{AF-E Node expressions} (\code{sh:path}, focus \code{sh:this},
  \code{sh:filterShape}, function application, \code{sh:intersection}/\code{union},
  constants) are \emph{value-producing} and map almost 1:1 onto $\pi$ plus
  literals plus function calls. They are the shared substrate for rule heads, AF
  targets, and computed values.
  \item \textbf{AF-T/AF-C/AF-CC/AF-F} (SPARQL targets, SPARQL constraints, custom
  components, SHACL functions) are escape-hatches: recognize a compilable subset
  back into the algebra, else evaluate opaquely; never silently wrong. JS
  features are out of scope with a clear diagnostic.
\end{itemize}

\paragraph{Invention worth flagging.} Class handling (C1) reduces what every
other SHACL engine treats as bespoke machinery (a type/subclass closure
subsystem) to \emph{one path expression}, so the same path planner, indexes, and
normalization apply to it for free.

% ===========================================================================
\section{Recursion semantics and semantic polarity}\label{sec:recursion}
% ===========================================================================

The paper is non-recursive; the W3C spec leaves recursion \emph{undefined}. Shape
references can be cyclic (\code{sh:node}/\code{sh:property}/\code{qualifiedValueShape}),
and the arena represents that natively. This section is load-bearing: every
optimization in \S\ref{sec:norm}+ is sound only relative to this choice.

\subsection{The decision}
\begin{enumerate}[leftmargin=1.6em]
  \item \textbf{Stratified} semantics over the shape/rule dependency graph, with
  a \emph{diagnostic} (not a guess) when a schema is not stratifiable.
  \item Within a stratum:
  \begin{itemize}
    \item \textbf{Validation} uses the \textbf{greatest fixpoint} ($\gfp$) ---
    coinductive ``no reachable counterexample.''
    \item \textbf{Inference} uses the \textbf{least fixpoint} ($\lfp$) --- only
    finitely-justified triples.
  \end{itemize}
  These run in separate phases (infer \emph{then} validate), so the directions
  never conflict.
\end{enumerate}

\subsection{Why stratified}
Pure positive recursion has clean fixpoints; \emph{negation through a cycle} does
not. The minimal paradox is
\[
S := \neg\,\exists p.\,S \qquad\text{over a $p$-self-loop,}
\]
which has no consistent 2-valued answer. A schema is \emph{stratifiable} iff its
shape-dependency graph has no cycle through a \emph{negative} edge. We evaluate
stratum by stratum (lower strata fully decided first). This is deterministic,
2-valued, optimizer-friendly, and --- crucially --- \textbf{the same machinery
serves AF rule inference} (Datalog with stratified negation), unifying recursive
validation and inference under one engine.

\subsection{Polarity must be computed \emph{semantically}}\label{sec:polarity}
This is the subtle part. Because the IR encodes $\forall\pi.\varphi$ as
$\exists^{\leq0}\pi.\neg\varphi$ (a \code{Count} with $\max{=}0$ over $\neg\varphi$),
a \emph{positive} SHACL constraint (\code{sh:node} $S$ in a property shape) looks
\emph{syntactically negative}. But two anti-monotone operators compose to
monotone, so the \emph{semantic} polarity is positive. The dependency analysis
tracks monotonicity, not surface $\neg$:
\begin{center}
\small
\begin{tabular}{@{}ll@{}}
\toprule
construct & polarity of the referenced shape \\
\midrule
\code{sh:node}/\code{sh:property}/\code{minCount}/$\exists^{\geq n}$ qualifier & \textbf{positive} (monotone) \\
\code{sh:not} & \textbf{negative} \\
\code{maxCount}/$\exists^{\leq n}$/qualified-max qualifier & \textbf{negative} (anti-monotone) \\
\code{closed}/\code{disjoint} & \textbf{negative} \\
\bottomrule
\end{tabular}
\end{center}
\begin{itemize}[leftmargin=1.4em]
  \item Because \code{Count\{min,max\}} carries \emph{one} qualifier governed by
  \emph{both} bounds, the analysis \textbf{un-fuses} it: the $\min$ part
  contributes a positive edge, the $\max$ part a negative edge; a qualifier under
  both yields a negative (non-monotone) edge.
  \item Each \code{DepEdge} carries a \code{Polarity} with sign $\pm1$. A cycle's
  net polarity is the \emph{product} of its edges' signs.
\end{itemize}

\subsection{The sign-balance stratification test}
\begin{principle}[Sign balance]
Condense strongly-connected components (Tarjan SCC) and order them by dependency
depth (dependees first); each SCC is one stratum. A recursive SCC is
\textbf{stratifiable} iff it is \emph{sign-balanced}: no cycle has net-negative
(odd) polarity. A sign-unbalanced SCC is recursion through genuine negation
(e.g.\ $S:=\neg\exists p.S$) and is reported.
\end{principle}
\begin{lstlisting}
struct Stratum { shapes: Vec<ShapeId>, recursive: bool, stratifiable: bool }
struct Stratification { strata: Vec<Stratum>, stratifiable: bool } // eval order: dependees first
\end{lstlisting}

\subsection{Why gfp for validation, lfp for inference}
$\lfp$ and $\gfp$ differ only on cyclic data (they coincide on DAGs). For
$\mathrm{conf}(v)\iff \mathrm{isPerson}(v)\wedge(\text{every \code{knows}-neighbour conforms})$
over a mutual \code{knows} cycle of two Persons:
\begin{itemize}[leftmargin=1.4em]
  \item $\lfp$ (build up from grounded bases): the cycle never grounds out
  $\Rightarrow$ both \textbf{fail} (inductive ``must bottom out'').
  \item $\gfp$ (whittle away provable violators): no reachable concrete violation
  $\Rightarrow$ both \textbf{conform} (coinductive ``no reachable
  counterexample'').
\end{itemize}
For validation the coinductive reading is what universal constraints usually
\emph{mean} and it avoids flagging legitimate cyclic data; hence $\gfp$ is the
default. Inference must be $\lfp$ (a rule fires only when its body is actually
satisfied; a fact cannot justify itself), which is also the standard semi-naive
fixpoint. Stratification supports either direction per positive stratum, so this
is \emph{not a one-way door}.

\paragraph{Engine realization.} The evaluator's ``assume conforming on a
back-edge'' cycle guard computes exactly the $\gfp$: on revisiting an active
$(\text{ShapeId},\text{node})$ pair it returns \code{true} (provisional), and ---
critically --- results that depend on a back-edge are \textbf{not memoized},
because their provisional truth may be valid only in the active call context.
\begin{lstlisting}
struct EvalState {
    memo:   HashMap<(ShapeId, Term), bool>,  // completed, context-free results
    active: HashSet<(ShapeId, Term)>,        // on the current recursion stack
}
// back-edge => EvalResult { holds: true, cacheable: false }  (the gfp choice)
\end{lstlisting}

% ===========================================================================
\section{Semantics-preserving algebraic transformations}\label{sec:norm}
% ===========================================================================

Normalization rebuilds the shape arena from the schema roots, hash-consing
structurally-identical nodes and applying sound Boolean/count/path/value-type
simplifications. \textbf{Soundness invariant}: all rules are per-node
truth-functional / relational, hence safe under the $\gfp$ validation semantics;
the W3C harness cross-checks
$\mathsf{validate}(\mathsf{normalize}(S),G)\equiv\mathsf{validate}(S,G)$ on every
core test (at the conforms + violation-focus level).

\subsection{Enabler: hash-consing (CSE) + compaction}
\begin{itemize}[leftmargin=1.4em]
  \item \textbf{CSE / hash-consing.} Structurally-identical arena nodes are
  deduplicated through a cons-table (interner). Acyclic nodes are interned
  bottom-up; recursive SCCs are rebuilt \emph{preserving sharing} but not
  collapsed, so cycles survive.
  \item \textbf{Arena compaction (free).} The rebuild interns only what the roots
  reach, so unreachable slots simply do not appear --- no orphans, no separate GC
  pass.
  \item CSE assigns stable ids that every later rule keys on (e.g.\
  complementation $\varphi\wedge\neg\varphi=\bot$ is detected by id equality).
\end{itemize}

\subsection{Boolean layer ($\varphi$)}
\begin{itemize}[leftmargin=1.4em]
  \item Flatten \code{And}/\code{Or}; $\neg\neg\varphi=\varphi$;
  $\code{And}([])=\top$, $\code{Or}([])=\bot$.
  \item $\top/\bot$ absorption: $\varphi\vee\top=\top$, $\varphi\wedge\bot=\bot$,
  $\varphi\vee\bot=\varphi$, $\varphi\wedge\top=\varphi$.
  \item Idempotent dedup $\varphi\wedge\varphi=\varphi$,
  $\varphi\vee\varphi=\varphi$ (sort children by id, dedup).
  \item Complementation $\varphi\wedge\neg\varphi=\bot$,
  $\varphi\vee\neg\varphi=\top$ (via CSE ids).
  \item Canonical child ordering (by id) so equal \code{And}/\code{Or} collapse.
  \item \textbf{NNF}: push $\neg$ inward with the count-flip laws
  \[
  \neg(\exists^{\geq n}\pi.\varphi)=\exists^{\leq n-1}\pi.\varphi,\qquad
  \neg(\exists^{\leq m}\pi.\varphi)=\exists^{\geq m+1}\pi.\varphi
  \]
  (flip bounds, keep qualifier positive). Negation is \emph{not} pushed into
  recursive nodes --- a recursive $\neg$ stays a literal, preserving
  stratifiability.
  \item Node-kind algebra: $\neg\,\kind(K)=\kind(\bar K)$ (or $\bot$ if empty);
  sibling intersection $\kind(K_1)\wedge\kind(K_2)=\kind(K_1\cap K_2)$ ($\bot$ if
  empty).
\end{itemize}

\subsection{Count layer ($\exists^{[\min..\max]}\pi.\varphi$)}
\begin{itemize}[leftmargin=1.4em]
  \item Trivial: $\min{=}0$ / unbounded $\Rightarrow\top$. Unsat bounds
  $\min>\max\Rightarrow\bot$.
  \item Qualifier collapse: $\exists^{\geq1}\pi.\bot=\bot$,
  $\exists^{\leq m}\pi.\bot=\top$.
  \item Same-$(\pi,\varphi)$ merge (conjunction context):
  $\exists^{\geq a}\wedge\exists^{\geq b}=\exists^{\geq\max(a,b)}$,
  $\exists^{\leq a}\wedge\exists^{\leq b}=\exists^{\leq\min(a,b)}$; a separate
  min/max count over the same $(\pi,\varphi)$ is fused into \emph{one} node so the
  path is evaluated once.
  \item $\id$-path collapse: $\exists^{\geq1}\id.\varphi=\varphi$,
  $\exists^{\leq0}\id.\varphi=\neg\varphi$.
  \item Empty-path zero law: $\exists^{\geq1}\,\code{Alt}([]).\varphi=\bot$.
\end{itemize}

\subsection{Path layer ($\pi$): Kleene + converse laws}
\begin{itemize}[leftmargin=1.4em]
  \item Flatten \code{Seq}/\code{Alt}; $\pi\cdot\id=\pi$; $(\pi^{-})^{-}=\pi$;
  $\id^{\Star}=\id$; $(\pi^{\Star})^{\Star}=\pi^{\Star}$ (smart constructors).
  \item \textbf{Converse push-down to canonical}:
  $(\pi_1\cdot\pi_2)^{-}=\pi_2^{-}\cdot\pi_1^{-}$, recursively, until $\code{Inverse}$
  wraps only a \code{Pred}. This means downstream code only ever sees
  $\code{Inverse}(\code{Pred}(q))$, never a compound inverse.
  \item \code{Alt} idempotence (dedup via a \code{HashSet}); empty-relation zero
  law.
  \item Star laws: $(\pi\cup\id)^{\Star}=\pi^{\Star}$ (strip $\id$ from an
  \code{Alt} before starring); $\pi^{\Star}\cdot\pi^{\Star}=\pi^{\Star}$ (merge
  adjacent equal stars in a \code{Seq}).
\end{itemize}

\subsection{Value-type layer ($\tau$)}
\begin{itemize}[leftmargin=1.4em]
  \item Flatten \code{And}, drop \code{Any}.
  \item Range/length \textbf{tightening}: fold same-family numeric/length bounds
  to the tighter pair; conjoined $\test(\tau)$ siblings are fused. Numeric
  ordering covers xsd numeric only --- incomparable bounds are kept, never
  mis-merged.
  \item Facet unsat $\Rightarrow\bot$: empty range, $\text{len.min}>\text{max}$,
  conflicting datatypes; $\test(\mathsf{any})\Rightarrow\top$.
  \item \textbf{Cross-facet} unsat: a known non-numeric datatype (e.g.\
  \code{xsd:string}) conjoined with a \code{NumericRange} is declared $\bot$.
  Unknown/custom datatypes are kept, never mis-declared.
\end{itemize}

\subsection{Schema layer}
\begin{itemize}[leftmargin=1.4em]
  \item Statement/selector dedup (identical $(\sel,\varphi)$ pairs after
  normalization).
  \item Selector canonicalization:
  $\code{HasPath}(\code{Pred}(q),\top)\Rightarrow\code{HasOut}(q)$ and
  $\code{HasPath}(\code{Pred}(q)^{-},\top)\Rightarrow\code{HasIn}(q)$.
\end{itemize}

\paragraph{Recursion caveat.} CSE across cycles is id-based and safe;
inlining/unfolding \emph{into} cycles is not attempted. Recursive SCCs get only
light rebuilding (no collapse). Normalization preserves stratifiability.
Reporting drift is accepted: normalization changes \code{@id}s/structure, so
per-constraint \emph{messages} can change; oracle equivalence is conforms +
violation-focus, not message-identical.

% ===========================================================================
\section{Static analysis and logical $\to$ physical planning}\label{sec:plan}
% ===========================================================================

\subsection{Backend-agnostic static analysis}
\nibox{The slicing / fingerprinting / context-footprint / shared-work-unit /
data-graph-summary subsystem below was built in the prior \code{shacl-core}
architecture and is \emph{not} present in the current \code{shifty-opt} /
\code{shifty-engine} crates. What \emph{is} implemented today is the
polarity-aware dependency graph and stratification (\S\ref{sec:recursion}), the
cost model and \code{And}/\code{Or} reordering, selector seeding
(\S\ref{sec:plan}), and per-query profiling/telemetry (query fingerprints in
\code{profile.rs}). Treat this subsection as a design target.}
Over a normalized program, the static analyzer can:
\begin{itemize}[leftmargin=1.4em]
  \item \textbf{Slice} irrelevant shapes/rules/constraints/targets/components from
  target roots (or explicit root shapes), recording per-item inclusion/exclusion
  reasons; recursive SCCs stay intact in the slice.
  \item \textbf{Classify the context footprint} of each shape/rule
  (node-local scalar vs.\ single-hop vs.\ bounded traversal vs.\
  shape-reference traversal vs.\ global/SPARQL).
  \item \textbf{Fingerprint} constraints/targets/rules with stable structural
  hashes, grouping structural duplicates into \emph{shared work units} (repeated
  SPARQL constraints, repeated component instantiations, repeated rule bodies).
  \item Emit static \textbf{cost hints} and separate \emph{validation-closure}
  vs.\ \emph{inference-closure} dependency classes (validation-only,
  inference-only, shared edges).
  \item Optionally fold in \textbf{data-graph statistics} (predicate
  cardinality, distinct subjects/objects, node-kind/datatype distributions,
  class/subclass counts) to mark empty/unsatisfiable target scans and obviously
  dead or vacuous constraints --- all advisory, never semantics-changing.
\end{itemize}

\subsection{Selector seeding (target-driven focus enumeration)}
The principal logical-planning move is to \emph{seed focus nodes from selectors}
instead of scanning all nodes:
\begin{lstlisting}
enum FocusSource {
    SubjectsOf(NamedNode),   // subjects of (., p, .)
    ObjectsOf(NamedNode),    // objects  of (., p, .)
    Node(Term),              // a single sh:targetNode
    PathToConst { path: Path, target: Term },  // class targets: backward seed
    ScanFilter  { path: Path, qualifier: ShapeId },  // general path selector
    Sparql(SparqlTarget),
}
\end{lstlisting}
\begin{principle}[Backward seeding of class targets]
For a class target $\exists^{\geq1}\pi.\test(c)$, instead of scanning every node
$v$ and testing $\exists\pi.\test(c)$ \emph{forward}, seed \emph{backward} from
the constant: the focus set is $\mathrm{pred}(c,\pi)=\{v\mid(v,c)\in\semG{\pi}\}$.
\end{principle}
This turns a whole-graph scan into a reverse path lookup from a single constant
($\approx 23\times$ over a scan on a synthetic class-target graph, in the current
build).

\subsection{Cost-based short-circuit ordering}
\code{And}/\code{Or} children are reordered cheapest-first using a static cost
model, so the short-circuiting evaluator rejects/accepts on cheap atoms before
walking expensive paths. The model (relative cost per arena node at one focus):
\[
\begin{array}{ll}
\top,\ \text{pending} \mapsto 0 &
\test(c),\test(\tau),\kind(K)\mapsto 1\\
\closedc(Q)\mapsto 4 &
\eqc/\disjc/\code{Lt}/\code{Le}\mapsto 2+\mathrm{cost}(\pi)\\
\code{Sparql}\mapsto 100 &
\text{back-edge (recursion)}\mapsto 50\\
\code{And}/\code{Or}\mapsto \textstyle\sum_i \mathrm{cost}(c_i) &
\code{Count}\mapsto \mathrm{cost}(\pi)\cdot(1+\mathrm{cost}(q))
\end{array}
\]
with path cost $\mathrm{cost}(\code{Pred})=1$,
$\mathrm{cost}(\pi^{-})=1+\mathrm{cost}(\pi)$,
$\mathrm{cost}(\code{Seq}/\code{Alt})=\sum\mathrm{cost}(\cdot)$, and
$\mathrm{cost}(\pi^{\Star})=10\cdot(1+\mathrm{cost}(\pi))$. Costs are memoized;
a recursion back-edge during cost computation is charged a flat 50.

\paragraph{Other plan-time work.} Implemented: common-subexpression sharing
across shapes (the CSE in \S\ref{sec:norm}), and analysis-informed deterministic
ordering of local/cheap before global/SPARQL constraints (the cost model charges
$\code{Sparql}{=}100$, sorting it last).
\NI{explicit $\test$/$\closedc$ filter pushdown as a plan rewrite, and
incremental / focus-delta scheduling so only affected focus nodes are re-checked
\emph{during validation}. (Predicate-delta scheduling \emph{is} implemented on
the \emph{inference} side --- \S\ref{sec:infer}.)}

% ===========================================================================
\section{Native SPARQL lowering --- the principal speed mechanism}\label{sec:sparql}
% ===========================================================================

Profiling the 223P/NIST workload shows that \textbf{validation-side SPARQL
evaluation dominates runtime}: the shapes lower to a small number of SPARQL
constraint leaves, $\approx99\%$ of sampled validation time is below SPARQL
constraint evaluation, and most self-time is Oxigraph quad iteration, range
checks, and property-path traversal --- concentrated in
\code{rdfs:subClassOf}$^{\Star}$, \code{rdf:type/rdfs:subClassOf}$^{\Star}$, and
\code{s223:cnx}$^{+}$/\code{contains}$^{+}$/\code{mapsTo}$^{+}$.

\paragraph{Decision.} Compile a useful Spargebra subset into a \emph{native
physical plan} over immutable, specialized indexes. Queries outside the subset
run through Spareval over the \emph{same} indexed dataset. A query is
\textbf{never} partly evaluated by both engines: capability analysis chooses one
executor before evaluation.

\subsection{Execution pipeline}
\[
\text{canonical SPARQL}\!\rightarrow\!\text{Spargebra AST}\!\xrightarrow{\text{static SHACL subst.}}\!\text{capability+demand}\!
\begin{cases}
\rightarrow \text{NativeQueryPlan}\rightarrow\text{native exec}\\
\rightarrow \text{PreparedSparqlQuery}\rightarrow\text{Spareval}
\end{cases}\!\rightarrow\text{FrozenIndexedDataset}
\]
Each canonical query is parsed once; the compilation cache key is
$(\text{canonical query},\ \text{static substitutions},\ \text{graph mode},\
\text{index-plan version})$, mapping to either a native plan or a prepared
fallback.

\subsection{SHACL prebinding as substitution, not a VALUES row}
SHACL prebinding is substitution \emph{throughout} the query, not an ordinary
initial binding. The compiler represents substituted values as plan operands:
\begin{itemize}[leftmargin=1.4em]
  \item \textbf{Static per constraint}: \code{\$PATH}, \code{\$currentShape},
  \code{\$shapesGraph}, custom-component parameters --- applied once before
  capability analysis. A complex \code{\$PATH} is substituted in
  \emph{property-path position}, not encoded as an RDF term.
  \item \textbf{Dynamic per invocation}: \code{\$this} --- native plans take a
  \code{FocusId}/\code{TermId} input column so one plan evaluates many focus nodes
  together.
\end{itemize}

\subsection{Capability analysis}
\code{analyze\_capability} walks the AST and returns \code{Native} or
\code{Fallback\{ first\_unsupported \}} (the first offending node, recorded for
telemetry). The native subset:
\begin{itemize}[leftmargin=1.4em]
  \item \code{SELECT}/\code{ASK}; BGPs; fixed named-graph patterns;
  \item joins, unions, projection, \code{DISTINCT};
  \item property paths: predicate, reverse, sequence, alternative, $*$, $+$, $?$;
  \item correlated \code{EXISTS}/\code{NOT EXISTS};
  \item boolean connectives, \code{BOUND}, safe equality, \code{STR},
  \code{STRSTARTS}; simple \code{BIND}; inline \code{VALUES}.
\end{itemize}
Fallback triggers on aggregates, \code{ORDER BY}, \code{LIMIT}/\code{OFFSET},
\code{OPTIONAL}, \code{MINUS}, \code{LATERAL}, \code{SERVICE}, variable graph
names, ordered comparisons, \code{IN}, and unsupported functions.

\subsection{The native physical plan}
The actual operator set (in \code{shifty-opt}) is a \emph{pipelined} plan: each
operator threads the previous operator's solutions through its \code{input},
which is sound because a BGP join only ever \emph{adds} bound variables.
\begin{lstlisting}
enum NativeOp {
    InputFocus,                          // one seed per focus node, binds $this
    Scan     { input, pattern: TripleScan }, // indexed nested-loop, threaded on input
    PathScan { input, scan: PathScan },      // arbitrary-length path (*/+/?)
    Union    { left, right },                 // over the same input
    Filter   { input, expr: ExprPlan },       // keep rows whose EBV is true
    Extend   { input, var, expr },            // BIND
    Project  { input, vars }, Distinct { input },
}
\end{lstlisting}
\begin{itemize}[leftmargin=1.4em]
  \item \textbf{Joins are pipelined scans, not a binary operator.} A multi-pattern
  BGP is lowered to a chain of \code{Scan}/\code{PathScan} nodes, each consuming
  the previous one's bound solutions (indexed nested-loop). Triple-pattern
  reordering by estimated cardinality is greedy from \code{DatasetStatistics}
  (\S\ref{sec:frozen}).
  \item \textbf{Batched focus.} Every solution carries a \code{FocusId}; one
  batched run keeps each focus node's results separate, so \code{DISTINCT},
  projection, and bucketing are per-focus --- as if the query had been run once
  per node with \code{\$this} pre-bound, but executed once.
  \item \textbf{\code{PathScan} binding modes (implemented).}
  bound-start+bound-end $\to$ membership probe; bound-start $\to$ forward closure;
  bound-end $\to$ reverse closure; open endpoints $\to$ relation scan over the
  node domain.
  \item \textbf{Correlated \code{EXISTS}/\code{NOT EXISTS}} are handled inside
  \code{Filter}'s \code{ExprPlan} (capability analysis admits them), not as a
  dedicated operator.
\end{itemize}
\nibox{The richer plan in \code{docs/05-sparql-execution.md} --- a dedicated
binary \code{Join \{ kind \}} with a \emph{hash}-join alternative, and a dedicated
\code{Exists} semi-join / anti-join operator with first-match early exit, and the
two-mode (\emph{Conforms} early-exit vs.\ \emph{Violations}) constraint executor
--- is design only. The current executor uses pipelined indexed nested-loop
scans and expression-level \code{EXISTS}.}

\subsection{Frozen indexed dataset}\label{sec:frozen}
Because inference runs to a fixed point \emph{before} validation, validation
receives an immutable graph. One dictionary-encoded dataset is built at that
boundary (actual struct):
\begin{lstlisting}
struct FrozenIndexedDataset {
    terms: TermDictionary,                 // Term <-> TermId (u32)
    default_graph: TripleIndex,            // SPO / POS / OSP sorted vecs
    named_graphs: HashMap<TermId, Vec<[TermId;3]>>, // small; SPO-sorted
    reach_cache: RefCell<ReachCache>,      // bounded per-start closure memo
    stats: DatasetStatistics,
}
struct TripleIndex { spo, pos, osp: Vec<[TermId; 3]> }  // three orders
struct DatasetStatistics {
    triple_count, distinct_subjects, distinct_objects: u64,
    predicate_cardinality: HashMap<TermId, u64>,
}
\end{lstlisting}
\begin{itemize}[leftmargin=1.4em]
  \item \textbf{Dictionary encoding.} \code{TermId = u32} is the
  \code{QueryableDataset::InternalTerm}, keeping each triple to 12 bytes. The
  dictionary is interior-mutable (\code{RefCell}) so query constants absent from
  the loaded triples are interned lazily under \code{\&self} (an unknown constant
  gets a fresh id that matches no stored triple --- exactly what a scan needs).
  \item \textbf{Three sorted immutable orders} (SPO, POS, OSP) cover the access
  patterns by range scan: $(s,p)\to o$ via SPO, $(p,o)\to s$ via POS,
  $o\to(s,p)$ via OSP, plus exact membership. Default and named graphs are
  separate.
  \item \textbf{Statistics} are computed at build time and \emph{used}: the native
  lowerer (\code{lower\_query\_with\_stats}) greedily reorders BGP triple patterns
  by estimated output cardinality from \code{predicate\_cardinality}. \PART{the
  statistics are predicate-cardinality + distinct subject/object counts only; no
  degree distributions or class/subclass counts.}
  \item \textbf{Shared backend.} \code{FrozenIndexedDataset} implements
  Spareval's \code{QueryableDataset}, so it backs \emph{both} native and fallback
  queries; the fallback gains fast immutable triple-pattern access even though its
  property paths are still Spareval-evaluated.
\end{itemize}

\subsection{Path demand and index planning}
Equivalent paths share a canonical \code{PathId} (the \code{path\_to\_string}
rendering, unique up to smart-constructor normalization). Demand is collected from
both the native algebra and Spargebra queries:
\begin{lstlisting}
struct PathDemand {
    path: PathId,
    forward_probes, reverse_probes, membership_probes, open_scans, expected_focuses: u64,
}
\end{lstlisting}
\PART{the \code{PathDemand} extractor is implemented as stage-1 \emph{static}
counts (per-occurrence in the AST); \code{expected\_focuses} is not yet populated
from dataset statistics.}
\nibox{The per-path \emph{strategy selection} below is design only. The current
engine has exactly two behaviors: base-index traversal (the \code{Traverse} row)
and the bounded per-start closure memoization described next. The \code{Memoized}
bitmap cache, \code{Materialized} forward/reverse CSR relations, and
\code{SccClosure} component-reachability indexes --- together with the named
path specializations and the budgeted, demand-driven promotion --- are \emph{not}
built.

\begin{center}\small
\begin{tabular}{@{}ll@{}}
\toprule
strategy & when \\
\midrule
\code{Traverse} & cheap or rarely used paths --- use base indexes directly \\
\code{Memoized} & cache forward/reverse result bitmaps per endpoint \\
\code{Materialized} & store the full relation in forward + reverse CSR form \\
\code{SccClosure} & condense a transitive graph into SCCs; store component reachability bitmaps \\
\bottomrule
\end{tabular}
\end{center}
Intended specializations: \code{rdfs:subClassOf}$^{\Star}$ $\to$ SCC condensation
+ component reachability; \code{rdf:type/rdfs:subClassOf}$^{\Star}$ $\to$
materialized instance$\leftrightarrow$class relations;
\code{s223:contains}$^{+}$/\code{mapsTo}$^{+}$/\code{cnx}$^{+}$ $\to$ memoized
endpoint closures, promoted to materialized only when demand and estimated size
justify it. Index selection would be budgeted: candidates ranked by avoided
traversal work per byte and admitted until a memory budget is exhausted, with
\code{Traverse} always the correct fallback.}

\paragraph{Closure memoization (implemented).} The executor implements bounded
per-start memoization for native $*$/$+$ closures:
\begin{itemize}[leftmargin=1.4em]
  \item Cache key $=(\text{start term},\ \text{compiled reachability step},\
  \text{closure kind},\ \text{graph selection})$.
  \item Endpoint sets are \emph{shared} between repeated probes via \code{Rc},
  capped at $10^6$ stored endpoint ids, and cleared if inference extends the
  dataset.
\end{itemize}

\subsection{Compiled reachability plans (path evaluation)}
\code{compile\_fwd}/\code{compile\_bwd} translate a \code{Path} into a
\code{ReachStep} tree \emph{once}, at executor-creation time, doing two things
that would otherwise recur per seed node per step:
\begin{enumerate}[leftmargin=1.6em]
  \item \textbf{Predicate interning} --- each \code{Pred(q)} is resolved to a
  \code{TermId} up front, eliminating a hash-map lookup per seed per step.
  \item \textbf{Direction folding} --- \code{Inverse} is absorbed into the step
  type (\code{FwdPred} vs.\ \code{BwdPred}) and \code{Seq} elements are reversed
  for backward traversal, so the evaluator never re-dispatches on direction.
\end{enumerate}
\begin{lstlisting}
enum ReachStep { Id, FwdPred(TermId), BwdPred(TermId), Seq(Vec<_>), Alt(Vec<_>), Star(Box<_>) }
// FwdPred(q): SPO range scan(node, q, _) -> objects
// BwdPred(q): POS range scan(_, q, node) -> subjects
\end{lstlisting}
The evaluator then walks the compiled tree with direct index calls; star/plus
closures are BFS loops shared between nested \code{Path::Star} and the top-level
\code{ClosureKind}.

\subsection{Correctness, fallback, and differential testing}
\begin{itemize}[leftmargin=1.4em]
  \item Oxigraph/Spareval remains the \textbf{semantic oracle}. A query uses
  fallback iff capability analysis cannot \emph{prove} every node/expression/graph
  op/result form supported.
  \item Native planning or execution errors are surfaced as engine errors ---
  \emph{never} silently downgraded to fallback after partial evaluation. Fallback
  is a \emph{planning} decision, keeping behavior deterministic.
  \item A differential mode (under \code{debug\_assertions}) runs native-capable
  queries through both engines and compares: \code{ASK} booleans; \code{SELECT}
  multisets (including unbound variables); RDF term/literal identity; SHACL
  violation values/paths; errors and unsupported prebindings.
\end{itemize}

% ===========================================================================
\section{SHACL-AF inference: semi-naive least fixpoint}\label{sec:infer}
% ===========================================================================

\begin{itemize}[leftmargin=1.4em]
  \item A rule fires on focus nodes for which all \code{sh:condition}s hold,
  producing triples from its head's node expressions. Inference is the
  \textbf{least fixpoint} (\S\ref{sec:recursion}).
  \item Rules run in ascending \code{sh:order} groups; output from a later group
  may reactivate an earlier group in the next pass.
  \item \textbf{Predicate-level delta scheduling.} A rule is rerun only when newly
  added triples can be observed by its graph reads; the first cut uses
  conservative predicate-delta scheduling at rule granularity. Native paths
  expose exact reads, while opaque/domain-sensitive constructs fall back to
  wildcard rescheduling.
  \item \textbf{Termination of the supported subset.} Triple rules only combine
  \emph{existing} terms; SPARQL \code{CONSTRUCT} results containing \emph{fresh}
  blank nodes are rejected. So no new domain elements are minted, bounding the
  derivable set. Function node expressions are reported as diagnostics, not
  silently skipped.
  \item Interaction model: infer-then-validate, with provenance for inferred
  triples. The same stratified framework drives both inference ($\lfp$) and
  validation ($\gfp$).
\end{itemize}

% ===========================================================================
\section{Symbolic repair: abduction over $\varphi$}\label{sec:repair}
% ===========================================================================

When $G\not\models S$, validation says \emph{that} a focus violates a shape and
which component failed. Repair answers the dual: \emph{the set of ways the data
graph could be edited so the focus satisfies the shape}. shifty ships this as a
\textbf{library, not an autonomous repairer} --- it computes and gates; a
\emph{driver} decides.

\subsection{Repair as the abductive dual of validation}
Validation is $G,v\conf\varphi$ by structural recursion on $\varphi$. Repair is:
\[
\Repair(\varphi,v)=\{\Delta G : (G\oplus\Delta G),v\conf\varphi\},
\]
where $\Delta G$ is a set of triple additions/deletions and $\oplus$ applies
them. The library \emph{describes} this set via the \emph{same} structural
recursion on $\varphi$ --- the small core (the \code{Shape} enum is $\sim$15
variants, already in NNF, a foldable DAG) pays off twice.
\begin{principle}[Compute from the algebra, not the report]
The W3C report walker treats \code{and}/\code{or}/\code{not}/\code{node} as opaque
units. Repair must drill in (to repair $\varphi_1\wedge\varphi_2$ it needs the
repair spaces of \emph{both} conjuncts), so it recurses over the arena, using the
report only to seed which statements failed at which focus nodes.
\end{principle}

\subsection{The witness / satisfaction-trace duality}
To describe repairs only for the parts of $\varphi$ that failed, the engine
computes two interlocked, mutually-recursive folds beside \code{explain}, sharing
the same \code{holds} oracle and back-edge guard:
\begin{itemize}[leftmargin=1.4em]
  \item \code{witness} (additive direction --- \emph{what to add}): the failed
  sub-DAG of $\varphi$, pruned to exactly what did not hold, carrying at each node
  the \code{ShapeId} and the structural gap (e.g.\ \code{CountLow\{path, have,
  min, qualifier\}}). It is the structured, lossless sibling of the human-facing
  \code{Reason}.
  \item \code{sat\_trace} (deletive direction --- \emph{what to delete}): a record
  of \emph{why $\varphi$ currently holds}, needed because a $\neg\varphi$ failure
  can only be repaired by \emph{breaking} $\varphi$.
\end{itemize}
The two folds return dual sum types. Each constructor mirrors a $\varphi$
constructor, pruned to the relevant sub-DAG; \code{PathSupport} records the
\emph{existing} edges that made a node a $\pi$-successor (the cut candidates).
\begin{lstlisting}
enum PathSupport { Empty, Edge(Triple), Chain(Vec<PathSupport>), Alt(Vec<PathSupport>) }

enum Witness {                                  // additive: what to ADD to make phi hold
    Atom       { shape, node, reached_by: Path, produced_by: Option<PathSupport> }, // test() leaf
    Relational { shape, node, kind: RelKind, lhs, rhs, offending },  // Eq/Disj/Lt/Le/UniqueLang
    Closed     { shape, node, offenders: Vec<(Pred, Term)> },
    Not        { shape, node, inner: Box<SatTrace> },          // crosses to the deletive side
    All        { shape, node, failed:   Vec<Witness> },        // conjunction: every child failed
    Any        { shape, node, branches: Vec<Witness> },        // disjunction: every branch failed
    CountLow   { shape, node, path, qualifier, have, min, sibling_qualifiers }, // under-satisfied (>= min)
    CountHigh  { shape, node, path, qualifier, matched, max, per_value }, // over-satisfied (<= max)
    Opaque     { shape, node },                                // sh:sparql -- no algebraic witness
}
enum SatTrace {                                 // deletive: what to DELETE (why phi holds)
    Irrefutable { shape },                                     // Top: nothing to cut
    Atom        { shape, node, reached_by: Path, produced_by: PathSupport },
    AllHeld     { shape, node, children:  Vec<SatTrace> },     // conjunction: all hold => break one
    AnyHeld     { shape, node, satisfied: Vec<SatTrace> },     // disjunction: these hold => break all
    CountHeld   { shape, node, path, qualifier, matches, min, max },
    NotHeld     { shape, node, inner_fails: Box<Witness> },    // crosses to the additive side
    Blocked     { shape, node, reason: BlockReason },          // closed/relational/sparql in scope
    Coinductive { shape, node },                               // gfp back-edge: no facts to delete
}
\end{lstlisting}
\noindent Two reading conventions for the examples (\S\ref{sec:repair-examples}):
the \code{qualifier} field is a \code{ShapeId}, rendered as the shape it points to
(e.g.\ \code{Top}, \code{test("admin")}); and \code{All}/\code{Any}/\code{Atom}
are deliberately reused across \code{Witness} and \code{RepairTree} --- they denote
the \emph{same} skeleton in the failure tree and the repair tree --- so
\code{RepairTree} nodes are written with their id (\code{All\#4}) and
\code{Witness}/\code{SatTrace} nodes without one.
\begin{principle}[Direction flips]
\code{witness} and \code{sat\_trace} are \textbf{exact complements}: for a given
$(\text{node},\text{id})$ exactly one returns \code{Some} (modulo the back-edge
convention). At a \code{Not} they cross --- $\code{witness}\Rightarrow
\code{sat\_trace}(c)\mapsto\code{Not}$ and
$\code{sat\_trace}\Rightarrow\code{witness}(c)\mapsto\code{NotHeld}$. \code{Count}
is the other flip point (self-dual: break the lower bound by \emph{deleting}
matches, the upper bound by \emph{adding} them).
\end{principle}
\textbf{Honest limit}: a $\gfp$ support reached on a back-edge is assumed true
without a grounded justification, so there is no finite set of facts to delete ---
deletion-direction repair is genuinely \emph{incomplete through positive
recursion} (reported as \code{Coinductive}).

\begin{principle}[Vacuous universals are obligations, not absences]
A \code{CountLow}\,$\exists^{\geq m}\pi.q$ that fails because $\pi$ has \emph{no}
values shares its conjunction with any universal $\forall\pi.\varphi$
(lowered to $\exists^{\leq0}\pi.\neg\varphi$, \S\ref{sec:gap}) on the
\emph{same} path. Those universals \emph{hold} when the value set is empty --- so
they never witness as failures --- yet \emph{every} value added to discharge the
\code{CountLow} is checked against them. Witnessing therefore records the inner
$\varphi$ of each such sibling in \code{CountLow.sibling\_qualifiers}, so synthesis
constrains the new value to $q$ \emph{and} all the $\varphi$'s rather than to the
deficient count's qualifier alone --- the difference between ``add some
\code{ex:p}'' and ``add an \code{ex:p} that is a \code{Temperature\_Sensor}.''
The collection is by \emph{conjunctive scope}, not lexical adjacency: \code{witness}
threads the universals down the tree, each \code{And} adding its $\forall\pi.\varphi$
children, so a count and a universal that sit in \emph{different} property shapes
(hence different \code{And} nodes) on the same path still meet. A disjunction does
\emph{not} extend the scope into its branches --- a $\forall\pi.\varphi$ under one
\code{sh:or} arm must not constrain a count under another --- which is exactly why
a flat ``all universals on $\pi$'' index would be unsound.
\end{principle}

\subsection{The template IR (\code{RepairTree})}
A parametric, \emph{inspectable} description of the repair space for one
violation: four constructs mirroring $\varphi$ (a repair tree is the skeleton of
a satisfaction proof), with typed holes and variadic blocks.
\begin{lstlisting}
struct NodeId(u32);   // a RepairTree node's stable address          (written #n)
struct Hole(u32);     // a typed placeholder a driver binds to a term (written ?n)
enum   Slot { Bound(Term), Open(Hole) }              // a triple position: fixed term or hole
struct TriplePattern { s: Slot, p: Slot, o: Slot }   // a triple whose slots may be holes
enum   Edit { Add(TriplePattern), Delete(TriplePattern) }  // each carries a driver-overridable Cost
enum HoleConstraint {                                // what a bound term must satisfy
    AnyNode, Fresh, Const(Term), Typed(ValueType), Kind(NodeKindSet),
    OneOf(BTreeSet<Term>), ConformsTo(ShapeId),       // value must satisfy a sub-shape
    ConformsToAll(Vec<ShapeId>),                      // ...or every one of several at once
}
enum RepairTree {
    Noop(NodeId), Blocked(NodeId, BlockReason),
    Edits  { id, edits: Vec<Edit>, holes: Vec<(Hole, HoleConstraint)> },
    All    { id, children: Vec<RepairTree> },  // satisfy every child
    Any    { id, children: Vec<RepairTree> },  // satisfy any one child
    Repeat { id, body: Box<RepairTree>, min: u64, max: Option<u64> },
}
\end{lstlisting}
\begin{itemize}[leftmargin=1.4em]
  \item \textbf{Holes and edits.} A \code{Hole} is a typed placeholder (written
  \code{?n}); the \code{HoleConstraint} paired with it in an \code{Edits} node
  states the admissible bindings --- a value type, a constant, \code{OneOf} a
  finite set, \code{ConformsTo} a sub-shape (the deferral point for
  \textbf{fuel}, below), or \code{ConformsToAll} a \emph{conjunction} of
  sub-shapes when one fresh value must discharge several obligations at once (a
  count qualifier plus its vacuous-universal siblings). An \code{Edits} node is a set of
  add/delete \code{TriplePattern}s whose \code{Open} slots are exactly those holes;
  a driver \emph{fills} a tree by choosing \code{Any}/\code{Repeat} options and
  binding every reachable hole.
  \item \textbf{Why \code{ConformsToAll}, and why no disjunctive dual.}
  \code{ConformsToAll} is not ``the $\wedge$ constraint'' for its own sake: it
  plugs a \emph{representational} gap. Its members are sibling \code{ShapeId}s
  collected independently during witnessing, and the arena is immutable during
  synthesis, so there is no slot standing for their conjunction to point a plain
  \code{ConformsTo} at. Disjunction never hits this gap: every shape-level
  alternative \emph{already} lowers to a single backing node --- \code{sh:or} and
  \code{sh:in} to an \code{Or}, \code{sh:xone} to the \code{Or} of its
  mutual-exclusion expansion $\bigvee_i(\varphi_i\wedge\bigwedge_{j\neq i}\neg\varphi_j)$ ---
  so a one-of-several-shapes hole is just \code{ConformsTo} of that node, and a
  failing \code{Or} is branched \emph{structurally} as an \code{Any} subtree, not
  carried on a hole. (Term-level disjunction is \code{OneOf}; the deletive
  \code{maxCount} hole is exactly that --- which existing values to drop --- so it
  needs no shape disjunction either.) The asymmetry is therefore real but benign:
  a \code{ConformsToOneOf} would have no source that produces un-backed disjuncts.
  \item \textbf{Blocked propagation invariant} (like the arena's smart
  constructors): \code{All} with any \code{Blocked} child $\Rightarrow$
  \code{Blocked}; \code{Any} drops \code{Blocked} children and is \code{Blocked}
  only if all are. So a returned subtree never hides a dead branch inside a
  satisfiable one; a top-level \code{Blocked} means ``no data repair exists in
  scope.''
  \item \code{BlockReason}: \code{OpaqueSparql}, \code{CannotMutateIdentity}
  (focus-scoped nodeKind/identity), \code{Coinductive}, \code{Unsupported}.
\end{itemize}

\subsection{Synthesis: three mutually-recursive folds}
\code{synthesize\_focus} is graph-free (all contact with $G$ happened during
witnessing) and decides nothing. It is:
\begin{itemize}[leftmargin=1.4em]
  \item $\Repair(\code{Witness})$ --- additive: make a \emph{failing existing}
  node hold. Walks the (already-pruned) witness; finite.
  \item $\Break(\code{SatTrace})$ --- deletive: falsify a \emph{holding existing}
  node. The De~Morgan mirror of $\Repair$.
  \item $\Build(\code{ShapeId},\code{Hole})$ --- additive, \emph{hypothetical}:
  constrain a \emph{not-yet-existing} node to satisfy a shape. Walks the
  \code{Shape} (everything must be constructed, so unlike $\Repair$/$\Break$ it can
  diverge on a recursive shape); this is the fold that \textbf{fuel} governs (below).
\end{itemize}

\paragraph{Fuel: bounding hypothetical construction.} $\Repair$ and $\Break$ are
structurally terminating because they walk an \emph{already-pruned} witness /
satisfaction trace --- finite objects produced by witnessing against a concrete
$G$. $\Build$ has no such bound: satisfying a \emph{not-yet-existing} value against
a shape can require satisfying \emph{another} fresh value against the \emph{same}
shape. A recursive qualifier (a \code{:Person} that must have a \code{:knows}-edge
to a node that is itself a \code{:Person}) turns naive structural expansion into an
infinite tower of minted nodes. \textbf{Fuel} is the depth budget that cuts this
off: a counter decremented each time the builder descends through a \code{Count}
qualifier into a fresh value, and when it reaches $0$ the still-recursive
obligation is emitted as a \code{ConformsTo(id)} hole --- an explicit ``some node
must satisfy shape \code{id} here'' --- rather than expanded further. The driver
then discharges that hole, typically by \emph{reuse} (binding it to a node already
conforming to \code{id}) instead of minting an unbounded chain. Fuel is therefore
the boundary between \emph{deferring} a recursive obligation and \emph{eagerly
materializing} it.

\PART{In the current engine \code{synthesize} carries \emph{no} fuel counter at
all: rather than expand recursive qualifiers to any depth, every \code{CountLow}
qualifier is lowered to a single \code{ConformsTo} hole at its first recursive step
(value-type leaves still get a precise \code{Typed}/\code{Kind} constraint), so the
fold bottoms out trivially at the witness/trace leaves --- it always
\emph{defers}. Fuel re-appears \emph{in the driver}: a driver that elects to
\emph{expand} a \code{ConformsTo} hole --- mint a fresh node and recursively build
it out against shape \code{id} (\code{RepairSession.repair\_node\_against}) ---
carries its own depth budget. The shipped interactive driver uses
\code{BUILD\_FUEL = 6} and halts with a ``build depth limit reached'' diagnostic at
$0$, falling back to reuse or manual authoring.}
Selected dispatch ($\Repair$):
\[
\begin{array}{ll}
\code{CountLow}\{\pi,q,\mathrm{have},\min\} & \mapsto \code{Repeat}\{\min{-}\mathrm{have},\ \code{build\_value}(\pi,\,q\wedge\textstyle\bigwedge_i\varphi_i)\}\\
\code{CountHigh (plain max)} & \mapsto \code{Repeat}\{\mathrm{have}{-}\max,\ \code{Delete over OneOf(matched)}\}\\
\code{CountHigh}\ (\forall\text{-encoding}) & \mapsto \code{All}(\text{per-value } \Repair(\text{inner}))\\
\code{Not}\{\text{inner}\} & \mapsto \Break(\text{inner})\\
\code{Atom (value-scoped)} & \mapsto \code{Any}([\text{delete offender, build a good value}])\\
\end{array}
\]
The \code{CountLow} qualifier $q\wedge\bigwedge_i\varphi_i$ is the count's own
qualifier conjoined with every sibling universal collected during witnessing
(above): $\top$ conjuncts drop out, a lone survivor becomes the precise leaf or
\code{ConformsTo} constraint, and several become a \code{ConformsToAll}. So when a
\code{sh:qualifiedValueShape} count and a sibling \code{sh:class} both govern the
empty path, the hole carries \emph{both} obligations, not the count's qualifier
alone --- the gate would reject a value missing either, and now the repair space
\emph{says so}.
\textbf{Path materialization} is the key structural move: when the governing path
$\pi$ is a \code{Seq}, adding ``a $\pi$-reachable value'' mints the interior nodes
too --- each interior step becomes a \code{Fresh} hole and a chained \code{Add},
so one \code{Count} over a multi-step path yields a deep repair block.
\code{Inverse} flips subject/object; \code{Alt} becomes an \code{Any};
\code{Star} is bounded. Its deletion-side dual, \code{PathSupport}, enumerates the
\emph{existing} triples that made \code{node} a $\pi$-successor (for \code{Seq}, an
OR over the chain --- cut any one edge).

\subsection{Two worked examples}\label{sec:repair-examples}
Both examples target one node and run the whole pipeline
\emph{shape\,+\,data $\to$ witness/trace $\to$ \code{RepairTree}}. The node ids
(\code{\#n}) and hole ids (\code{?n}) shown are exactly those synthesis assigns
(both examples are regression tests in \code{synthesize.rs}). A single-constraint
node shape is \emph{not} wrapped in an \code{And}, so the witness trees below are
bare leaves rather than singleton conjunctions.

\paragraph{Example 1 --- additive: an under-satisfied count.} A minimum count with
too few values is the canonical \emph{additive} repair.
\begin{lstlisting}[language={}]
ex:S a sh:NodeShape ; sh:targetNode ex:x ;
     sh:property [ sh:path ex:p ; sh:minCount 2 ] .   # shape: >=2 values on ex:p
ex:x ex:p ex:y .                                       # data:  only one present
\end{lstlisting}
Lowering yields $\exists^{\geq 2}\,\code{ex:p}.\top$. Witnessing \code{ex:x} counts
one matching value where two are required, so the witness is a single
\code{CountLow} leaf:
\begin{lstlisting}[language={}]
CountLow{ node: ex:x, path: ex:p, qualifier: Top, have: 1, min: 2 }
\end{lstlisting}
$\Repair$ turns the deficit $\min-\mathrm{have}=1$ into a \code{Repeat} whose body
\emph{materializes} one fresh \code{ex:p}-value; the value hole \code{?0} carries
the qualifier's constraint --- here \code{AnyNode}, because the qualifier is $\top$:
\begin{lstlisting}[language={}]
Repeat#1 { min: 1, max: None
  body: Edits#0 { add( ex:x  ex:p  ?0 ),  hole ?0 : AnyNode } }
\end{lstlisting}
A driver fills this by choosing the \code{Repeat} count ($1$) and binding \code{?0}
--- to an existing node (reuse) or a freshly minted one --- giving the single add
\code{ex:x ex:p <v>}. Had the qualifier been a value type (e.g.\
\code{sh:qualifiedValueShape [ sh:datatype xsd:integer ]}), \code{?0} would instead
carry \code{Typed(xsd:integer)}; and a \code{Seq} path would mint the interior hops
as \code{Fresh} holes (path materialization).

\paragraph{Example 2 --- the witness/trace duality through \code{sh:not}.} The
illuminating case is a violation fixable only by \emph{deletion}, which is exactly
where the \code{SatTrace} surfaces. ``This account must \emph{not} be flagged
banned'':
\begin{lstlisting}[language={}]
ex:S a sh:NodeShape ; sh:targetNode ex:acct ;
     sh:not [ sh:path ex:status ; sh:hasValue "banned" ] .  # shape: must NOT hold
ex:acct ex:status "banned" .                                # data:  but it does
\end{lstlisting}
Lowering yields
$\neg\,\exists^{\geq 1}\,\code{ex:status}.\test(\text{``banned''})$. The inner
$\exists$ \emph{holds}, so the $\neg$ \emph{fails} --- and witnessing therefore
crosses onto the satisfaction side: a \code{Not} witness embeds the \code{SatTrace}
that records \emph{why} the inner shape holds. This \textbf{is} the duality --- a
negation failure is an inner \emph{success} that must be dismantled:
\begin{lstlisting}[language={}]
Not{ node: ex:acct,
  inner: CountHeld{ path: ex:status, qualifier: test("banned"), min: 1,
           matches: [ ("banned", Atom{ produced_by: edge(ex:acct ex:status "banned") }) ] } }
\end{lstlisting}
Synthesis flips folds at the \code{Not}: $\Repair(\code{Not})=\Break(\text{inner})$.
$\Break(\code{CountHeld})$ with $\min>0$ drops the count below the minimum by
breaking every match, and $\Break(\code{Atom})$ cuts the edge that produced the
value --- a concrete deletion, no holes:
\begin{lstlisting}[language={}]
Edits#0 { delete( ex:acct  ex:status  "banned" ) }
\end{lstlisting}
The dual subtlety: a \emph{bare} \code{sh:not} over \code{minCount} alone
(qualifier $\top$) yields an \code{Irrefutable} match trace, and \code{break\_}
returns \code{Blocked} --- there is no value-type leaf to cut, only the path edge,
which this cut does not delete to falsify a bare $\exists$.

\paragraph{Example 3 --- both directions in one violation.} A node shape with
several constraints produces an \code{All} witness, and when one of those
constraints is a \code{sh:not} over a compound shape, that single violation report
carries a \emph{witness above the \code{Not}} and a \emph{satisfaction trace below
it} --- the two halves of the duality side by side. Policy: an account must carry
an email \emph{and} must not be an admin without MFA.
\begin{lstlisting}[language={}]
ex:S a sh:NodeShape ; sh:targetNode ex:acct ;
     sh:property [ sh:path ex:email ; sh:minCount 1 ] ;   # (A) >=1 email
     sh:not [ sh:and (                                     # (B) NOT (admin AND no MFA)
        [ sh:path ex:role ; sh:hasValue "admin" ]
        [ sh:path ex:mfa  ; sh:maxCount 0 ] ) ] .
ex:acct ex:role "admin" .       # an admin, with no MFA and no email -> fails (A) and (B)
\end{lstlisting}
The node shape lowers to
$\exists^{\geq1}\code{ex:email}.\top \;\wedge\;
\neg\big(\exists^{\geq1}\code{ex:role}.\test(\text{``admin''}) \wedge
\exists^{\leq0}\code{ex:mfa}.\top\big)$. Both conjuncts fail (no email; and the
account \emph{is} an un-MFA'd admin, so the negated inner shape holds), giving an
\code{All} witness whose second child is a \code{Not} embedding the \code{SatTrace}
for the inner \code{And}:
\begin{lstlisting}[language={}]
All{ node: ex:acct                                              <-- Witness (additive) world
  CountLow{ path: ex:email, qualifier: Top, have: 0, min: 1 }              # (A)
  Not{ node: ex:acct                                                       # (B): inner holds
    inner: AllHeld{                                             <-- SatTrace (deletive) world
      CountHeld{ path: ex:role, qualifier: test("admin"), min: 1,
        matches: [ ("admin", Atom{ produced_by: edge(ex:acct ex:role "admin") }) ] },
      CountHeld{ path: ex:mfa,  qualifier: Top, max: 0, matches: [] } } } }
\end{lstlisting}
Synthesis recurses over the \code{All}, producing one repair per failed conjunct
and conjoining them --- here an \emph{additive} \code{Repeat} for (A) and a
\emph{deletive} \code{Edits} for (B), in a single \code{All} repair:
\begin{lstlisting}[language={}]
All#4 {
  Repeat#1 { min: 1, max: None                                             # (A): add an email
    body: Edits#0 { add( ex:acct  ex:email  ?0 ),  hole ?0 : AnyNode } },
  Edits#2 { delete( ex:acct  ex:role  "admin" ) } }                        # (B): drop admin role
\end{lstlisting}
Two things to read off the (B) branch. First, $\Break(\code{AllHeld})$ is an
\code{Any} over its children's breaks --- breaking \emph{any one} conjunct
falsifies the \code{And}, hence un-admins the account. Second, the \code{ex:mfa}
\code{maxCount}-$0$ conjunct can only be falsified by \emph{adding} an MFA value,
which the deletive fold will not do, so it synthesizes \code{Blocked} (consuming
node id \code{\#3}); the enclosing \code{Any} then drops it (the
\code{Blocked}-propagation invariant of \S\ref{sec:repair}), which is why the ids
jump \code{\#2}$\to$\code{\#4} and why the surviving fix is ``drop the admin role''
rather than ``add MFA.'' The repair is sound regardless of which the gate accepts.

\subsection{The gate (whole-graph, returned not acted on)}
A repair is sound only if it does not trade one violation for another, so the gate
re-validates over the whole graph:
\[
\code{gate}(G,S,\Delta G)=\big(\underbrace{\text{fixed}}_{\text{removed}},\ \underbrace{\text{introduced}}_{\text{new}},\ \underbrace{\text{remaining}}_{\text{still unfixed}}\big)
= \mathrm{violations}(G\oplus\Delta G,S)\ominus\mathrm{violations}(G,S).
\]
A delta is \emph{sound} iff \code{introduced} is empty; sound + non-empty
\code{fixed} is \emph{progress}. The gate reuses \code{validate} unchanged, so an
\emph{affected-set} re-validation (only nodes $\Delta G$ can touch, from the
dependency/demand analysis) can replace it later with identical semantics.

\subsection{Rendering to a self-hosted SHACL sidecar}
\nibox{\code{render} $\to$ \code{RepairProgram} (the parameterized-RDF view, the
\code{urn:shifty:hole\#} parameter nodes, the self-hosted SHACL sidecar over the
holes, and DNF flattening) is design only --- it is not in the current
\code{shifty-repair} / \code{shifty-engine} crates. The witness/\code{sat\_trace}
folds, the \code{RepairTree}/\code{Plan}/\code{instantiate} IR, \code{synthesize},
and \code{gate} \emph{are} implemented; drivers consume those directly.}
For graph-matching/display drivers, \code{render} restates the tree as a
parameterized RDF graph:
\begin{itemize}[leftmargin=1.4em]
  \item \textbf{Holes as parameter nodes}: each \code{Hole} renders to
  \code{urn:shifty:hole\#<id>} (the BuildingMOTIF \code{urn:\_\_\_param\_\_\_\#}
  analogue); a \code{TriplePattern} renders with bound slots as terms, open slots
  as the hole IRI.
  \item \textbf{Constraints as a self-hosted sidecar}: a \code{HoleConstraint}
  \emph{is} a shape, so it renders as a node shape on the hole IRI \emph{in the
  engine's own algebra} (\code{Const}$\to\test(c)$, \code{Typed}$\to$value-type,
  \code{Kind}$\to$\code{sh:nodeKind}, \code{OneOf}$\to$\code{sh:in},
  \code{ConformsTo}$\to$\code{sh:node}, \code{ConformsToAll}$\to$ several
  \code{sh:node}). No second constraint language; the
  sidecar is a \code{Schema} over the holes, checked by the \emph{same} validator
  the repair targets.
  \item \textbf{Flattening to DNF}: \code{Any}$\to$alternatives;
  \code{All}$\to$merge add-graphs/sidecars/delete-sets (contradictions prune the
  branch); \code{Repeat}$\to$count-parameterized block.
\end{itemize}

\subsection{Reference drivers}
\nibox{Of the four drivers below, only \emph{enumeration}-style filling is backed
by code today: \code{Hole.candidates()} (reuse-first candidates from the data
graph) is implemented, and two simple Python example drivers ship
(\code{python/examples/repair.py}, take-first-option; and
\code{repair\_interactive.py}, a numbered gated menu). The \emph{monomorphism},
\emph{ASP}, and \emph{LLM} drivers --- including the entire ASP lowering below
(\code{viol}/\code{conf}/\code{reach} Datalog, generate--define--test,
\code{\#minimize}) --- are specified in \code{docs/07-repair-drivers.md} as
worked examples, not implemented in the tree.}
The library is solver-agnostic; drivers bring data and policy and run the loop
(witness $\to$ synthesize $\to$ fill \code{Plan} $\to$ instantiate $\to$ gate
$\to$ accept/re-witness). Termination is the driver's, but easy: require each
accepted $\Delta G$ to shrink the violation set (\code{gate} reports this).
\begin{center}
\small
\begin{tabular}{@{}llll@{}}
\toprule
driver & fills holes by & strength & joint? \\
\midrule
Monomorphism & matching existing subgraphs & reuse / structural overlap & no (greedy) \\
Enumeration & brute force over finite domains & simplicity, IR oracle & no (greedy) \\
LLM & model proposal & semantic / infinite holes & no (greedy) \\
ASP & constraint solving & minimality, cascades & \textbf{yes} \\
\bottomrule
\end{tabular}
\end{center}

\paragraph{The ASP driver --- joint solving + minimality.} The key observation is
that \emph{lowering $\varphi$ to a logic program is mechanical because the
recursion semantics is already stratified negation}, and the engine already runs
it as least-fixpoint forward chaining with a computed stratum order: $\varphi$
\emph{is} a logic program; lowering reuses the same strata.
\begin{itemize}[leftmargin=1.4em]
  \item \textbf{Generate} (template-bounded): choose hole bindings, repeat-instance
  counts, and deletions \emph{restricted to what the \code{RepairTree} proposes}
  --- not ``any triple in the universe'' --- keeping the search finite.
  \item \textbf{Define}: post-repair \code{holds/3}, the path relations
  \code{reach/3}, and \code{viol/2} (= the witness) as Datalog over \code{holds}.
  \item \textbf{Test}: \code{:- not conf(root, focus).} for every root shape in the
  focus group (the conjoined integrity constraints \emph{are} the joint solve).
  \item \textbf{Optimize}: \code{\#minimize} over chosen edits with reuse$<$mint
  weights.
\end{itemize}
\begin{principle}[Lower \code{viol}, not \code{conf}]
Validation is $\gfp$ but ASP stable models are $\lfp$. Lowering \code{viol} (the
finitely-justified failure --- exactly the witness) and defining \code{conf} as
its stratified negation makes the reading $\lfp$/finitely-justified, which is the
\emph{correct} semantics for repair: a repair must actually \emph{construct}
supporting structure, not lean on a coinductive cycle. The divergence is contained
because the driver only runs on genuine $\gfp$ witnesses and every answer set
still passes the $\gfp$ \code{gate}.
\end{principle}
The non-logical leaf checks ASP cannot express (\code{sh:datatype}, numeric range,
\code{sh:pattern}, \code{sh:nodeKind}, \code{sh:in}, \code{Lt}/\code{Le}/\code{Eq})
are \textbf{precomputed by the engine} and emitted as facts (\code{sat(Leaf,X)},
\code{lt(X,Y)}, \dots), reusing the same \code{value\_type\_holds}/\code{compare\_terms}
the validator uses. ASP decides structure/reuse/counts; novel infinite literals
are left to the enumeration/LLM drivers.

% ===========================================================================
\section{Catalog of invention (what was pulled together vs.\ created)}
% ===========================================================================

\paragraph{Pulled together (adapted from existing theory).}
\begin{itemize}[leftmargin=1.4em]
  \item The $\pi$/$\varphi$ core and its denotational semantics, from the
  \emph{Common Foundations} paper --- but \textbf{specialized to RDF} (D0/D1
  re-unification) rather than common graphs.
  \item Stratified negation / $\lfp$ Datalog evaluation (classical), reused
  verbatim for both AF inference and ASP repair lowering.
  \item Dictionary-encoded, sorted, multi-order triple indexes (standard triple
  store engineering), specialized to an \emph{immutable} post-inference snapshot.
\end{itemize}

\paragraph{Created / non-obvious for this engine.}
\begin{enumerate}[leftmargin=1.6em]
  \item \textbf{Class-as-path (C1).} Eliminating the class/subclass subsystem by
  expressing \code{sh:class}/\code{targetClass} as a single
  \code{rdf:type}$\cdot$\code{rdfs:subClassOf}$^{\Star}$ path, so the path planner
  and indexes serve it for free.
  \item \textbf{Fused \code{Count}} as the single counting primitive subsuming
  min/max/qualified cardinality, \code{node}/\code{property} nesting, and
  $\forall/\exists$ --- with deliberate \emph{un-fusing} for polarity analysis.
  \item \textbf{Semantic polarity \& sign-balance stratification.} Computing
  monotonicity through the $\exists^{\leq0}\pi.\neg\varphi$ encoding of $\forall$
  (so a syntactically-negative IR node is correctly positive), and testing
  stratifiability as net-sign balance of SCC cycles.
  \item \textbf{$\gfp$-validation / $\lfp$-inference split} under one stratified
  framework, with a non-cacheable back-edge guard that realizes the $\gfp$ exactly.
  \item \textbf{Backward seeding of class targets} (reverse path lookup from the
  class constant) instead of node scans.
  \item \textbf{Native SPARQL subset over a frozen indexed dataset}: capability
  analysis as a whole-query gate, substitution-based prebinding,
  \code{FocusId}-batched execution, compiled direction-folded reachability plans,
  bounded shared-\code{Rc} closure memoization, and stats-driven greedy BGP
  reordering --- plus differential testing against Spareval as the oracle.
  \PART{the \emph{budgeted multi-strategy per-path index catalog} (Memoized /
  Materialized / SccClosure) is design only; see \S\ref{sec:sparql}.}
  \item \textbf{Symbolic repair as abduction over $\varphi$}: the
  witness/satisfaction-trace duality with provable add/delete flips at
  \code{Not}/\code{Count}, the fuel-governed \code{build} for hypothetical nodes
  (recursion deferred via \code{ConformsTo} holes; \S\ref{sec:repair}),
  path materialization / \code{PathSupport} as additive/deletive duals, and the
  whole-graph \code{gate} returned-not-acted-on. \NI{the self-hosted SHACL sidecar
  rendering of hole constraints.}
  \item \NI{\textbf{ASP repair lowering that reuses the validation strata},
  lowering \code{viol} (not \code{conf}) to get $\lfp$/finitely-justified repair
  semantics while still gating against the $\gfp$ validator --- specified, not
  built.}
\end{enumerate}

% ===========================================================================
\section*{Open problems / future work}
\addcontentsline{toc}{section}{Open problems / future work}
% ===========================================================================
\begin{itemize}[leftmargin=1.4em]
  \item \textbf{Demand-driven index promotion} (Memoized $\to$ Materialized $\to$
  SccClosure) from \emph{observed} demand is designed but only partially
  implemented (bounded per-start memoization exists; promotion is planned).
  \item \textbf{Affected-set re-validation} for the gate and for incremental
  focus-delta scheduling.
  \item \textbf{Repair completeness through positive recursion}
  (\code{Coinductive} limit) and finer \code{sh:equals} set-reconciliation.
  \item \textbf{Layer 7 compilation/JIT}: staged closures then optional Cranelift
  codegen for hot validators; vectorized batch evaluation; parallel strata.
  \item \textbf{Provenance side-table}: per-\code{sh:ValidationResult} provenance
  (the IR is pure, so reporting is currently focus-level for the algebra path).
\end{itemize}

% ===========================================================================
\section*{References (informal)}
\addcontentsline{toc}{section}{References}
% ===========================================================================
\begin{itemize}[leftmargin=1.4em]
  \item S.\ Ahmetaj et al., \emph{Common Foundations for SHACL, ShEx, and
  PG-Schema}, arXiv:2502.01295. (Source of the $\pi$/$\varphi$ core.)
  \item W3C, \emph{Shapes Constraint Language (SHACL)}, \url{https://www.w3.org/TR/shacl/}.
  \item W3C, \emph{SHACL Advanced Features}, \url{https://www.w3.org/TR/shacl-af/}.
  \item W3C data-shapes test suite, \url{https://w3c.github.io/data-shapes/data-shapes-test-suite/}.
  \item Project design docs: \code{docs/00-formalism.md} through
  \code{docs/07-repair-drivers.md}, \code{docs/static-analysis-plan.md}.
\end{itemize}

\end{document}
